\documentclass[10pt,letterpaper]{article}

% ── PLoS Computational Biology formatting ──
\usepackage[margin=1in]{geometry}
\usepackage{amsmath,amssymb,amsfonts}
\usepackage{graphicx}
\usepackage{booktabs}
% algorithm/algorithmic not available; use verbatim for pseudocode
\usepackage{fancyvrb}
\usepackage{hyperref}
\usepackage{natbib}
\usepackage{xcolor}
\usepackage{listings}
\usepackage{subcaption}
\usepackage{multirow}
\usepackage{array}
\usepackage{bm}
\usepackage{mathtools}
\usepackage{lineno}
\linenumbers

\graphicspath{{../benchmarks/}}

% ── Code listing style ──
\lstset{
  basicstyle=\ttfamily\footnotesize,
  keywordstyle=\color{blue!70!black}\bfseries,
  commentstyle=\color{green!50!black}\itshape,
  stringstyle=\color{red!60!black},
  breaklines=true,
  frame=single,
  captionpos=b,
  numbers=left,
  numberstyle=\tiny\color{gray},
  language=C++,
  morekeywords={__global__,__device__,__shared__,__forceinline__,
    __shfl_down_sync,__shfl_sync,__expf,__logf,__popcll,
    atomicAdd,cudaMalloc,cudaMemcpy,cudaFree,
    float,double,uint64_t,int,constexpr,template}
}

% ── Macros ──
\newcommand{\Ne}{N_\mathrm{e}}
\newcommand{\tmax}{t_\mathrm{max}}
\newcommand{\tmid}{\bar{t}}
\newcommand{\E}{\mathbb{E}}
\newcommand{\R}{\mathbb{R}}
\newcommand{\calB}{\mathcal{B}}
\newcommand{\calO}{\mathcal{O}}

\title{%
  \texttt{tmrca.cu}: GPU-Accelerated Pairwise Coalescence Time Estimation\\
  at Biobank Scale via Batched Hidden Markov Models\\
  on the Sequentially Markov Coalescent%
}

\author{%
  % Author list placeholder
  Author Name$^{1,\ast}$\\[4pt]
  {\small $^1$~Affiliation}\\
  {\small $^\ast$~Corresponding author: \texttt{email@example.com}}
}

\date{}

\begin{document}

\maketitle

% ════════════════════════════════════════════════════════════════
\begin{abstract}
Inferring pairwise times to the most recent common ancestor (TMRCA)
along the genome is fundamental to population genetics, yet existing
methods scale poorly to the sample sizes of modern biobanks
($n > 10^4$ haplotypes, $S > 10^6$ segregating sites).  We present
\texttt{tmrca.cu}, a CUDA library that formulates TMRCA inference as a
batched hidden Markov model (HMM) on the Sequentially Markov
Coalescent (SMC) and executes it entirely on GPU hardware.  The system
implements a four-tier architecture: (1)~instant windowed divergence
via bitpacked prefix scans, (2)~PELT changepoint segmentation,
(3)~full posterior inference with per-site credible intervals via the
SMC~HMM forward--backward algorithm, and (4)~adaptive coalescent
prior estimation via expectation--maximization.  Five phases of kernel
optimization---emission precomputation with fast-math intrinsics,
single-precision arithmetic, multi-pair thread blocks, fused summary
extraction, and persistent GPU memory contexts---yield a measured
throughput of $93 \times 10^6$~site$\cdot$pairs/s on an NVIDIA A100,
representing a $\mathbf{10.2\times}$ speedup over the na\"ive
double-precision baseline.  At this rate, all $\binom{1000}{2} =
499{,}500$ pairwise TMRCA profiles across human chromosome~1
($\sim\!4.5 \times 10^6$ SNPs) complete in approximately 7~hours on a
single GPU.  We validate numerical accuracy against a pure-NumPy
reference implementation across four demographic scenarios (constant
$\Ne$, bottleneck, exponential growth, structured populations),
demonstrating per-site Pearson correlations $r > 0.8$ with true
coalescence times, well-calibrated 95\% credible intervals (coverage
$\sim\!0.93$), and faithful recovery of non-trivial coalescent prior
distributions.  The software is implemented in CUDA~C++17 with
pybind11 Python bindings and is freely available under an open-source
license.
\end{abstract}

\noindent\textbf{Keywords:} coalescence, TMRCA, GPU, CUDA, hidden Markov model,
sequentially Markov coalescent, population genetics, biobank

% ════════════════════════════════════════════════════════════════
\section{Introduction}
\label{sec:intro}

The ancestral recombination graph (ARG) encodes the complete
genealogical history of a sample of chromosomes, and its marginal
trees determine pairwise coalescence times $T_{ij}(s)$ at each
genomic position~$s$.  These times are of direct biological interest:
they govern patterns of identity-by-descent (IBD), inform
demographic inference~\citep{Li2011,Schiffels2014,Terhorst2017},
underlie methods for detecting natural selection~\citep{Stern2019},
and provide the natural distance metric for phylogeographic
analysis.

Several methods infer local genealogies or pairwise TMRCAs from
sequence data.  The pairwise sequentially Markov coalescent
(PSMC)~\citep{Li2011} and its multi-sample extensions
(MSMC~\citep{Schiffels2014}, MSMC2~\citep{Schiffels2020}) model
$T_{ij}(s)$ as a hidden Markov model along the genome, with
discrete time bins as hidden states and observed
heterozygosity/homozygosity as emissions.  More recent methods infer
full ARGs (e.g., \texttt{tsinfer}~\citep{Kelleher2019},
\texttt{Relate}~\citep{Speidel2019}, \texttt{ARG-Needle}~\citep{Zhang2023}),
but their computational cost grows
superlinearly in sample size, limiting practical application to
$n \lesssim 10^4$.

A fundamental bottleneck is that the number of pairwise HMM
computations scales as $\binom{n}{2}$, and each requires $\calO(SK)$
operations where $S$ is the number of segregating sites and $K$ is the
number of time bins.  For a biobank cohort with $n = 10^3$ haplotypes
and $S = 4.5 \times 10^6$ SNPs (human chromosome~1), this amounts to
$\sim\!2.25 \times 10^{12}$ site$\cdot$pair evaluations---far beyond
the reach of CPU-based implementations.

Graphics processing units (GPUs) are well-suited to this problem:
the forward--backward algorithm for each pair is independent, and the
per-site computation involves a fixed number of floating-point
operations across $K$ time bins (typically $K = 32$, matching the GPU
warp size).  However, realizing this potential requires careful
attention to memory layout, arithmetic precision, warp-level
communication, and the elimination of unnecessary data transfers
between host and device.

In this paper, we describe \texttt{tmrca.cu}, a CUDA library that
implements GPU-accelerated pairwise TMRCA inference via batched HMMs
on the SMC\@.  Our contributions are:

\begin{enumerate}
  \item A \textbf{four-tier inference architecture} spanning instant
    divergence estimates (Tier~1), changepoint segmentation (Tier~2),
    full posterior inference (Tier~3), and adaptive prior learning
    (Tier~4), with ultrametric tree constraints enforced via
    expectation propagation (EP).

  \item A \textbf{highly optimized CUDA kernel} for batched
    forward--backward that exploits emission precomputation,
    single-precision fast-math intrinsics (\texttt{\_\_expf},
    \texttt{\_\_logf}), multi-pair thread blocks, fused on-the-fly
    posterior summary extraction, and persistent GPU memory contexts,
    achieving $93$M~site$\cdot$pairs/s on an A100 ($10.2\times$ over
    baseline).

  \item \textbf{Comprehensive validation} against a pure-NumPy
    reference implementation and msprime~\citep{Baumdicker2022}
    ground-truth genealogies under four demographic models, with
    quantitative accuracy metrics (correlation, RMSE, CI coverage,
    prior recovery).
\end{enumerate}


% ════════════════════════════════════════════════════════════════
\section{Model}
\label{sec:model}

\subsection{Sequentially Markov Coalescent}
\label{sec:smc}

We follow the pairwise SMC formulation~\citep{Li2011,McVean2005}.
For a pair of haplotypes $(i,j)$ observed at $S$ segregating sites with
physical positions $x_1 < x_2 < \cdots < x_S$, the coalescence time
$T_{ij}(s)$ is modeled as a Markov chain along the genome.  The
continuous coalescence time is discretized into $K$ bins with
boundaries
\begin{equation}
  t_k = \tmax \left(\frac{k}{K}\right)^2, \quad k = 0, 1, \ldots, K,
  \label{eq:time_grid}
\end{equation}
where $\tmax = 10\,\Ne$ by default.  The quadratic spacing
concentrates resolution at recent times where data are most
informative~\citep{Schiffels2014}.  Bin midpoints are
$\tmid_k = (t_k + t_{k+1})/2$.

\subsection{Coalescent prior}
\label{sec:prior}

Under a constant effective population size $\Ne$ and the standard
diploid coalescent, the probability that the coalescence time falls in
bin~$k$ is
\begin{equation}
  q_k = \exp\!\left(-\frac{t_k}{2\Ne}\right)
      - \exp\!\left(-\frac{t_{k+1}}{2\Ne}\right),
  \label{eq:coal_prior}
\end{equation}
normalized so that $\sum_{k=0}^{K-1} q_k = 1$.  This prior serves as
the initial state distribution and as the target of recombination
events (Section~\ref{sec:transitions}).

\subsection{Emission model}
\label{sec:emissions}

\paragraph{Site emission.}
At segregating site~$s$, let $d_s = g_{i,s} \oplus g_{j,s} \in
\{0,1\}$ denote the XOR of the two haploid genotypes (1~if the pair
differs, 0~otherwise).  Under the infinite-sites mutation model with
per-base-pair per-generation rate~$\mu$,
\begin{equation}
  P(d_s = 1 \mid T = \tmid_k) = 1 - e^{-2\mu\,\tmid_k}, \qquad
  P(d_s = 0 \mid T = \tmid_k) = e^{-2\mu\,\tmid_k}.
  \label{eq:site_emission}
\end{equation}
The factor of~2 accounts for the two lineages accumulating mutations
independently.

\paragraph{Gap emission.}
Between consecutive segregating sites $s$ and $s+1$ separated by
$L_s = x_{s+1} - x_s - 1$ base pairs of unobserved sequence, the
probability that neither lineage carries a mutation in the gap is
\begin{equation}
  P_\mathrm{gap}(s, k) = e^{-2\mu\,\tmid_k\,L_s}.
  \label{eq:gap_emission}
\end{equation}
This term is essential for correct likelihood computation when sites
are unevenly spaced~\citep{Li2011}.

\subsection{Transition model}
\label{sec:transitions}

Between sites $s$ and $s+1$, the recombination distance is
$r_s = \rho\,(x_{s+1} - x_s)$ where $\rho$ is the per-base-pair
per-generation recombination rate.  Under the SMC, recombination
events occur at rate proportional to the coalescence time:
\begin{equation}
  A_{kl}(r_s) = \begin{cases}
    e^{-r_s\,\tmid_k} + (1 - e^{-r_s\,\tmid_k})\,q_l & \text{if } k = l,\\
    (1 - e^{-r_s\,\tmid_k})\,q_l                     & \text{if } k \neq l.
  \end{cases}
  \label{eq:transition}
\end{equation}
The first term represents the ``stay'' component (no recombination),
and the second represents recombination followed by recoalescence
according to the prior~$q$.  This factorization is critical for
efficient GPU implementation: the recombination mass
$m_k = (1 - e^{-r_s\,\tmid_k})\,\alpha_{s-1}(k)$ can be summed
across all $K$ bins via a single warp reduction, then redistributed
as $q_k \cdot \sum_{l} m_l$ (Section~\ref{sec:kernel_forward}).


% ════════════════════════════════════════════════════════════════
\section{Methods}
\label{sec:methods}

\subsection{System architecture}
\label{sec:architecture}

\texttt{tmrca.cu} implements a four-tier inference pipeline with
increasing accuracy and computational cost (Fig.~\ref{fig:architecture}):

\begin{description}
  \item[Tier~1: Instant divergence.] Windowed fraction of differing
    sites, computed via bitpacked XOR prefix scans.  Runs in
    milliseconds; provides a rough $\hat{T}_{ij}(s) =
    \pi_{ij}(s,W)/(2\mu)$.

  \item[Tier~2: Changepoint segmentation.] Piecewise-constant TMRCA
    per pair via the Pruned Exact Linear Time (PELT)
    algorithm~\citep{Killick2012} with a Poisson cost function.  Each
    pair is processed by one warp on the GPU.

  \item[Tier~3: SMC HMM posterior.] Full posterior marginals
    $\gamma_{ij}(s,k) = P(T_{ij}(s) = \tmid_k \mid \text{data})$ via
    the forward--backward algorithm, yielding per-site means and 95\%
    credible intervals.

  \item[Tier~4: Adaptive prior via EM.] Iterative refinement of the
    coalescent prior~$q$ from the data, using the EM algorithm with
    fused posterior accumulation on the GPU\@.
\end{description}

Additionally, an \textbf{expectation propagation} (EP) loop couples
the per-pair HMMs (Tier~3) with an across-pair ultrametric tree
constraint, enforcing genealogical consistency at each site.

\subsection{Genotype representation: bitpacking}
\label{sec:bitpack}

Haploid genotypes $G \in \{0,1\}^{n \times S}$ are packed into
$\texttt{uint64\_t}$ words, with each word encoding 64~consecutive
sites in LSB-first order:
\begin{equation}
  W_{i,w} = \sum_{b=0}^{63} g_{i,\,64w+b} \cdot 2^b.
  \label{eq:bitpack}
\end{equation}
This achieves $8\times$ compression over \texttt{uint8} storage and
enables hardware-accelerated XOR and popcount operations.  The packing
kernel launches a $(n_\mathrm{words}/4) \times (n/256)$ grid with
$(4, 256)$ thread blocks, achieving fully coalesced memory access.

\subsection{Pairwise XOR prefix scan (Tier~1)}
\label{sec:prefix_scan}

For each pair $(i,j)$, the cumulative number of differing sites up to
position~$s$ is computed via a two-phase GPU algorithm:

\begin{enumerate}
  \item \textbf{Per-word XOR + popcount:} Each thread in a 32-thread
    warp processes a stripe of words via
    $c_w = \texttt{\_\_popcll}(W_{i,w} \oplus W_{j,w})$.

  \item \textbf{Warp-level inclusive prefix scan:} Using
    \texttt{\_\_shfl\_up\_sync}, the partial sums are accumulated
    across the warp with $\calO(\log_2 32) = 5$ shuffle steps.
\end{enumerate}

Windowed divergence is then a simple subtraction:
$\pi_{ij}(s,W) = [C_{ij}(s+W) - C_{ij}(s-W)] / (2W)$.

\subsection{PELT changepoint detection (Tier~2)}
\label{sec:pelt}

For each pair, the XOR bit sequence $\mathbf{d}_{ij}$ is modeled as a
Poisson process with piecewise-constant rate (proportional to TMRCA).
The PELT algorithm~\citep{Killick2012} finds the optimal segmentation
minimizing
\begin{equation}
  \sum_{\text{segments}}\!\mathrm{Cost}_\mathrm{Poisson}(\text{segment})
  + |\calB|\,\beta,
  \label{eq:pelt_objective}
\end{equation}
where $\beta = \log S$ (BIC penalty) and the Poisson negative
log-likelihood for a segment $[a,b)$ with $C$~mutations over
$L$~base pairs is
\begin{equation}
  \mathrm{Cost}(a,b) = \begin{cases}
    -C\log(C/L) + C & \text{if } C > 0,\\
    0                & \text{if } C = 0.
  \end{cases}
  \label{eq:pelt_cost}
\end{equation}
The MLE coalescence time per segment is $\hat{T} = C / (2\mu L)$.

Each pair is assigned one 32-thread warp.  The pruning set
$\mathcal{R}$ (candidate changepoints) is maintained in shared memory;
its expected size is $\calO(1)$ per step~\citep{Killick2012}, yielding
$\calO(S)$ amortized complexity.  Warp-level reductions via
\texttt{\_\_shfl\_down\_sync} find the minimum cost across candidates
at each step.

\subsection{Batched HMM forward--backward (Tier~3)}
\label{sec:hmm}

The core of the system is a GPU kernel implementing the
forward--backward algorithm for batched pairs.

\subsubsection{Forward pass}
\label{sec:kernel_forward}

\paragraph{Initialization.}
At site $s = 0$:
\begin{equation}
  \alpha_0(k) = q_k \cdot P(d_0 \mid \tmid_k),
  \label{eq:forward_init}
\end{equation}
normalized so that $\sum_k \alpha_0(k) = 1$, with
$\ell \leftarrow \log\!\sum_k \alpha_0(k)$ initializing the
log-likelihood accumulator.

\paragraph{Recursion.}
For $s = 1, \ldots, S-1$:
\begin{align}
  \alpha_s(k) &= \Bigl[
    \underbrace{e^{-r_s\,\tmid_k} \cdot \alpha_{s-1}(k)}_{\text{stay}}
    + \underbrace{q_k \cdot \textstyle\sum_{l=0}^{K-1}
      (1 - e^{-r_s\,\tmid_l})\,\alpha_{s-1}(l)}_{\text{recombine}}
  \Bigr] \nonumber\\
  &\qquad\times P_\mathrm{gap}(s,k) \cdot P(d_s \mid \tmid_k),
  \label{eq:forward_recursion}
\end{align}
followed by normalization
$\alpha_s(k) \leftarrow \alpha_s(k) / Z_s$ where
$Z_s = \sum_{k'} \alpha_s(k')$, and
$\ell \leftarrow \ell + \log Z_s$.

The recombination mass $\sum_l (1 - e^{-r_s\,\tmid_l})\,\alpha_{s-1}(l)$
is a scalar computed by a single warp/block reduction across the $K$
threads (one per time bin).  This avoids the na\"ive $\calO(K^2)$
matrix--vector product.

\paragraph{Storage.}
Forward values $\alpha_s(k)$ are written to a scratch buffer
$\texttt{alpha\_store}[s \cdot K + k]$ for use during the backward
pass.

\subsubsection{Backward pass}
\label{sec:kernel_backward}

\paragraph{Initialization.}
$\beta_{S-1}(k) = 1$ for all~$k$.

\paragraph{Recursion.}
For $s = S-2, \ldots, 0$:
\begin{align}
  \beta_s(k) &= e^{-r_{s+1}\,\tmid_k} \cdot
    \underbrace{\beta_{s+1}(k) \cdot P(d_{s+1} \mid \tmid_k) \cdot P_\mathrm{gap}(s\!+\!1,k)}_{b_k}
  \nonumber\\
  &\quad + (1 - e^{-r_{s+1}\,\tmid_k}) \cdot
    \underbrace{\textstyle\sum_{l=0}^{K-1} q_l \cdot b_l}_{\text{total}},
  \label{eq:backward_recursion}
\end{align}
normalized as $\beta_s(k) \leftarrow \beta_s(k) / \sum_{k'} \beta_s(k')$.

\subsubsection{Posterior marginals and summaries}
\label{sec:posterior}

The posterior is computed site-by-site during the backward sweep:
\begin{equation}
  \gamma_s(k) = \frac{\alpha_s(k)\,\beta_s(k)}%
    {\sum_{k'}\alpha_s(k')\,\beta_s(k')}.
  \label{eq:posterior}
\end{equation}

Three posterior summaries are extracted \emph{on the fly} during the
backward pass (Section~\ref{sec:fused_summary}):
\begin{align}
  \hat{T}_s &= \sum_{k=0}^{K-1} \gamma_s(k)\,\tmid_k
    &\text{(posterior mean)}, \label{eq:mean}\\
  T_s^{(\mathrm{lo})} &: \sum_{k=0}^{k_\mathrm{lo}} \gamma_s(k) \geq 0.025
    &\text{(2.5\textsuperscript{th} percentile)}, \label{eq:ci_lower}\\
  T_s^{(\mathrm{hi})} &: \sum_{k=0}^{k_\mathrm{hi}} \gamma_s(k) \geq 0.975
    &\text{(97.5\textsuperscript{th} percentile)}. \label{eq:ci_upper}
\end{align}

\subsection{Adaptive prior via EM (Tier~4)}
\label{sec:em}

The coalescent prior $q$ may not match the true (possibly
non-constant) demography.  We refine it iteratively:

\begin{enumerate}
  \item Run batched forward--backward with current prior~$q^{(t)}$.
  \item Accumulate the empirical posterior:
    $\hat{q}_k^{(t)} = \frac{1}{|\text{pairs}| \cdot S}
    \sum_{p,s} \gamma_{p,s}(k)$.
  \item Blend:
    $q_k^{(t+1)} = (1-\alpha)\,q_k^{(t)} + \alpha\,\hat{q}_k^{(t)}$,
    then normalize.
\end{enumerate}

The accumulation in step~2 is performed \emph{within the kernel} via
per-thread local accumulators and a final \texttt{atomicAdd} to a
length-$K$ buffer (Section~\ref{sec:fused_em}), completely eliminating
the need for a separate aggregation pass over the
$n_\mathrm{pairs} \times S \times K$ posterior array.

\subsection{Expectation propagation with ultrametric constraints}
\label{sec:ep}

At any genomic site, the pairwise TMRCAs for $m$~haplotypes must be
consistent with a binary tree (the ultrametric constraint): for any
triple $(i,j,k)$, the two largest pairwise times are equal.  We
enforce this via an EP loop that alternates between:

\begin{enumerate}
  \item \textbf{HMM factor} (per-pair, parallel): Run
    forward--backward with per-site messages
    $\mathbf{m}_{ij}(s) \in \R_+^K$ modifying the prior at each site.

  \item \textbf{Ultrametric factor} (per-site, parallel): Project the
    marginal posteriors $\{\gamma_{ij}(s)\}$ to the nearest
    ultrametric tree via agglomerative clustering that maximizes
    $\sum_{(i,j) \in \text{merge}} \log \gamma_{ij}(s, \tau)$ subject
    to merge-time monotonicity.

  \item \textbf{Message update:} Compute EP factors as the ratio of
    the projected posterior to the HMM marginal, with damping
    $\alpha \in (0,1)$ and clamping to $[0.1, 10]$ for numerical
    stability.
\end{enumerate}

Convergence is assessed by the maximum absolute change in messages:
$\Delta = \max_{p,s,k} |m^{(t)}_{p,s,k} - m^{(t-1)}_{p,s,k}|$.
Empirically, 3--5 iterations suffice.


% ════════════════════════════════════════════════════════════════
\section{GPU Kernel Design and Optimization}
\label{sec:gpu}

The HMM forward--backward kernel is the computational bottleneck.
We describe five optimization phases applied cumulatively, each
contributing a measurable speedup (Table~\ref{tab:speedup}).

\subsection{Kernel template parameterization}
\label{sec:template}

The kernel is a C++ template with three compile-time parameters:
\begin{lstlisting}[language=C++,caption={Kernel signature.}]
template<int K_TPL, int PPB, int MODE>
__global__ void hmm_forward_backward_kernel(
    const uint64_t* __restrict__ packed, int n_words,
    const double* __restrict__ positions, int S,
    const double* __restrict__ mu,
    const double* __restrict__ cum_rho,
    const double* __restrict__ time_midpoints,
    const double* __restrict__ coal_prior,
    const float* __restrict__ messages,
    const int* __restrict__ pair_i_arr,
    const int* __restrict__ pair_j_arr,
    int n_pairs,
    float* __restrict__ gamma_out,
    double* __restrict__ log_lik_out,
    float* __restrict__ tmrca_mean_out,
    float* __restrict__ tmrca_lower_out,
    float* __restrict__ tmrca_upper_out,
    double* __restrict__ q_accum_out);
\end{lstlisting}

\noindent where:
\begin{itemize}
  \item \texttt{K\_TPL} $\in \{32, 64, 128\}$: number of time bins
    (matches or is a multiple of the warp size).
  \item \texttt{PPB} (pairs per block): number of independent pairs
    processed by each thread block.  Dispatched as
    $(K\!=\!32, \mathrm{PPB}\!=\!4)$,
    $(K\!=\!64, \mathrm{PPB}\!=\!2)$,
    $(K\!=\!128, \mathrm{PPB}\!=\!1)$, yielding 128~threads/block
    in all cases.
  \item \texttt{MODE} $\in \{0, 1, 2\}$: output mode
    (\textsc{Full\_Gamma}, \textsc{Summary\_Only},
    \textsc{Summary\_And\_EM}).
\end{itemize}

The grid is launched with $\lceil n_\mathrm{pairs} / \mathrm{PPB} \rceil$
blocks.  Each thread is identified by its pair lane
$p = \lfloor \texttt{threadIdx.x} / K \rfloor$ and time-bin index
$k = \texttt{threadIdx.x} \bmod K$.

\subsection{Phase~1: Emission precomputation and fast-math intrinsics}
\label{sec:phase1}

The emission probability $e^{-2\mu\,\tmid_k}$ depends only on the
time bin~$k$, not on the site or pair.  In the baseline kernel, this
exponential is recomputed 8~times per site (forward + backward,
match + mismatch, site + gap).  We precompute three registers per
thread at kernel entry:
\begin{lstlisting}[language=C++,caption={Emission precomputation.}]
const float mu0 = (float)mu[0];
const float t_k = (float)time_midpoints[k];
const float emit_match_k   = __expf(-2.0f * mu0 * t_k);
const float emit_mismatch_k = 1.0f - emit_match_k;
const float gap_base_k      = -2.0f * mu0 * t_k;
\end{lstlisting}

During the forward and backward loops, site emissions become a
branch-free register lookup:
\begin{lstlisting}[language=C++]
float emit = d ? emit_mismatch_k : emit_match_k;
\end{lstlisting}

Gap emissions require only one fast exponential per site (conditional
on $L_s > 0$):
\begin{lstlisting}[language=C++]
float gap_emit = (gap_bp > 0.0f) ? __expf(gap_base_k * gap_bp) : 1.0f;
\end{lstlisting}

Remaining \texttt{exp()} calls (transition: $e^{-r_s\,\tmid_k}$) are
replaced with \texttt{\_\_expf()}, and \texttt{log()} for the
normalizer is replaced with \texttt{\_\_logf()}.  The hardware
fast-math unit on the A100 executes these in $\sim\!1$~cycle versus
$\sim\!10$~cycles for the libm equivalents.

\textbf{Estimated speedup:} $3\text{--}4\times$ (eliminates 4 of 8
double-precision exponentials per site per thread).

\subsection{Phase~2: Single-precision arithmetic}
\label{sec:phase2}

The NVIDIA A100 delivers 19.5~TFLOPS in FP64 versus 39.0~TFLOPS in
FP32 (a $2\!:\!1$ ratio).  Since the forward--backward algorithm
normalizes probabilities at every site, all intermediate values
remain well-conditioned in $[0, 1]$.  We convert:
\begin{itemize}
  \item Forward $\alpha_s(k)$, backward $\beta_s(k)$, posterior
    $\gamma_s(k)$: \texttt{float}.
  \item All warp/block reductions: \texttt{float} overloads.
  \item Shared memory: \texttt{float} (halving footprint).
\end{itemize}

Only the log-likelihood accumulator remains in \texttt{double} to
avoid catastrophic cancellation over $\sim\!10^6$ summands, and the EM
accumulator uses \texttt{double} \texttt{atomicAdd} to preserve
accuracy across $\sim\!10^9$ contributions.

\textbf{Estimated speedup:} $\sim\!2\times$ from the FP32/FP64
throughput ratio.

\subsection{Phase~3: Multi-pair thread blocks}
\label{sec:phase3}

With $K = 32$, one pair requires exactly one warp (32~threads).  A
single-warp block underutilizes the A100's SM, which can schedule up
to 64~warps.  We pack $\mathrm{PPB} = 4$ pairs into each 128-thread
block:

\begin{center}
\begin{tabular}{cccl}
\toprule
$K$ & PPB & Block size & Communication \\
\midrule
32  & 4   & 128        & Warp shuffles only, zero \texttt{\_\_syncthreads} \\
64  & 2   & 128        & Shared memory + \texttt{\_\_syncthreads} \\
128 & 1   & 128        & Shared memory + \texttt{\_\_syncthreads} \\
\bottomrule
\end{tabular}
\end{center}

Each warp independently processes its own pair, sharing only L1 cache
for read-only data (positions, $\mu$, cumulative $\rho$, time
midpoints, prior).  For $K = 32$, all reductions are pure warp shuffles
with no synchronization barriers.

\textbf{Estimated speedup:} $1.3\text{--}1.5\times$ from improved SM
occupancy and L1 cache sharing.

\subsection{Phase~4: Fused summary extraction and EM accumulation}
\label{sec:fused_summary}

\subsubsection{Eliminating the gamma array transfer}

In \textsc{Full\_Gamma} mode, the kernel writes a
$n_\mathrm{pairs} \times S \times K$ array of \texttt{float} values,
which must be copied to the host (D2H) for summary extraction.  For
1000~pairs at $S = 50{,}000$ and $K = 32$, this array occupies
$\sim\!6.4$~GB---a prohibitive bandwidth cost.

In \textsc{Summary\_Only} mode, posterior summaries are computed
\emph{during the backward pass}, immediately after
Eq.~\eqref{eq:posterior}:

\paragraph{Mean.}
Each thread~$k$ computes $\gamma_s(k) \cdot \tmid_k$; a warp
reduction yields $\hat{T}_s$; thread~0 writes the result.

\paragraph{95\% credible interval.}
For $K = 32$, each thread broadcasts its $\gamma_s(k)$ to all lanes
via \texttt{\_\_shfl\_sync}, and a sequential CDF scan identifies the
2.5\textsuperscript{th} and 97.5\textsuperscript{th} percentiles:
\begin{lstlisting}[language=C++,caption={CDF scan via warp shuffle ($K\!=\!32$).}]
float cum = 0.0f;
for (int kk = 0; kk < 32; kk++) {
    float gk = __shfl_sync(0xFFFFFFFF, gamma_val, kk);
    cum += gk;
    if (!lower_set && cum >= 0.025f) {
        lower_val = time_midpoints[kk];
        lower_set = true;
    }
    if (cum >= 0.975f) {
        upper_val = time_midpoints[kk];
        break;
    }
}
\end{lstlisting}

For $K > 32$, the gamma values are written to the alpha scratch buffer
and thread~0 performs a sequential scan.

\subsubsection{Fused EM accumulation}
\label{sec:fused_em}

In \textsc{Summary\_And\_EM} mode, each thread additionally maintains
a local accumulator:
\begin{lstlisting}[language=C++]
float q_accum = 0.0f;
// During backward loop, for each site:
q_accum += gamma_val;
// After loop completion:
atomicAdd(&q_accum_out[k], (double)q_accum);
\end{lstlisting}

This yields the un-normalized empirical prior
$\hat{q}_k \propto \sum_{p,s} \gamma_{p,s}(k)$ in a single kernel
launch, eliminating the separate \texttt{aggregate\_posteriors\_gpu}
kernel that would otherwise require a full re-read of the posterior
array.

\textbf{Estimated speedup:} $\sim\!1.5\times$ from eliminating the
6.4~GB D2H transfer and the aggregation pass.

\subsection{Phase~5: Persistent GPU context and memory management}
\label{sec:persistent}

The \texttt{HMMContext} C++ class (exposed to Python via pybind11)
holds all static data on the GPU across multiple \texttt{run\_batch()}
calls:

\begin{itemize}
  \item Bitpacked genotypes (\texttt{uint64\_t*}).
  \item Positions, mutation rates, cumulative recombination distances
    (\texttt{double*}).
  \item Time midpoints and coalescent prior (\texttt{double*}).
\end{itemize}

Per-batch scratch buffers (alpha scratch, log-likelihood, summary
arrays, pair index arrays) are allocated once and reused.  An
auto-chunking mechanism queries available GPU memory via
\texttt{cudaMemGetInfo}, reserves 512~MB headroom, and splits large
pair sets into chunks that fit:
\begin{equation}
  n_\mathrm{chunk} = \left\lfloor
    \frac{\text{free memory} - 512\,\text{MB}}%
         {S \cdot K \cdot 4 + S \cdot 4 \cdot 3 + 8 + 8}
  \right\rfloor.
  \label{eq:chunking}
\end{equation}

\textbf{Estimated speedup:} $\sim\!1.25\times$ from eliminating
per-call host-to-device transfers and bitpacking.

\subsection{Warp-level and block-level reductions}
\label{sec:reductions}

The kernel relies on two reduction primitives:

\paragraph{Warp reduction (float).}
A tree-reduction using 5~shuffle steps:
\begin{lstlisting}[language=C++]
__device__ __forceinline__
float warp_reduce_sum_f(float val) {
    for (int offset = 16; offset > 0; offset >>= 1)
        val += __shfl_down_sync(0xFFFFFFFF, val, offset);
    return __shfl_sync(0xFFFFFFFF, val, 0);  // broadcast
}
\end{lstlisting}

\paragraph{Block reduction (templated on $K$).}
For $K = 32$, this reduces to a single warp reduction.  For
$K \in \{64, 128\}$, a two-stage cascade is used: (1)~warp-level
reduction within each of $K/32$ warps, (2)~lane-0 writes to shared
memory, (3)~warp-0 reduces across warp results.  The \texttt{k\_local}
parameter (the thread's bin index within its pair, not
\texttt{threadIdx.x}) is essential for correct warp/lane identification
in multi-pair blocks.

\subsection{XOR bit caching}
\label{sec:xor_cache}

Each pair caches the current 64-bit XOR word across consecutive sites:
\begin{lstlisting}[language=C++]
int cur_word_idx = -1;
uint64_t cur_xor_word = 0;
// ...
int w = s >> 6;
if (w != cur_word_idx) {
    cur_xor_word = packed[hi*n_words + w]
                 ^ packed[hj*n_words + w];
    cur_word_idx = w;
}
int d = (int)((cur_xor_word >> (s & 63)) & 1ULL);
\end{lstlisting}

This amortizes 64~bit extractions to 2~global memory reads (one per
haplotype), improving memory efficiency by $\sim\!32\times$ for the
genotype data.

\subsection{Loop unrolling}
\label{sec:unroll}

Both forward and backward site loops carry the pragma
\texttt{\#pragma unroll 2}, instructing the CUDA compiler to unroll by
a factor of~2.  This reduces loop overhead, improves instruction-level
parallelism, and enables better register scheduling, at the cost of
modest additional register pressure.


% ════════════════════════════════════════════════════════════════
\section{Ancillary GPU Kernels}
\label{sec:ancillary}

\subsection{Site frequency spectrum}
\label{sec:sfs}

Two kernel variants are provided:
\begin{itemize}
  \item \textbf{Simple} ($n \leq 256$): One thread per site counts
    the allele frequency by iterating over all $n$~haplotypes, using
    \texttt{atomicAdd} to update the SFS histogram.
  \item \textbf{Fast} ($n > 256$): A $(32, 8)$ thread block per word
    uses warp-level popcount reduction, processing $n/32$ haplotypes
    per thread and avoiding serialization on the histogram.
\end{itemize}

\subsection{Ultrametric projection}
\label{sec:ultrametric_kernel}

The per-site ultrametric projection is implemented as a
single-block-per-site kernel using shared memory for the $m \times m$
pair posteriors, cluster assignments, and merge tracking.  The
agglomerative clustering loop performs $m-1$ merges, each requiring a
scan over $\calO(m^2 K)$ candidates with block-wide reduction to find
the optimal merge.  The merge time must be monotonically non-decreasing
(constrained by \texttt{cluster\_max\_t}), ensuring tree
consistency.

\subsection{Convergence monitoring}
\label{sec:convergence_kernel}

A two-stage reduction kernel computes
$\Delta = \max_{p,s,k} |m^{(t)} - m^{(t-1)}|$ over the message
array.  Stage~1 computes per-block maxima (256~threads, 2~elements
each); stage~2 reduces the block maxima to a single scalar.  This
avoids a costly D2H transfer of the full message array.


% ════════════════════════════════════════════════════════════════
\section{Software Implementation}
\label{sec:software}

\subsection{Build system}

The project uses \texttt{pixi}~\citep{prefix_pixi} for reproducible
dependency management and CMake + Ninja for CUDA/C++ compilation.  The
CUDA toolkit is pinned to version ${\geq}\,12.0, {<}\,12.3$
(matching the A100 driver), and the target architecture is
\texttt{sm\_80}.  Separable compilation
(\texttt{CUDA\_SEPARABLE\_COMPILATION ON}) is enabled to allow
independent compilation of 8~kernel files.

\subsection{Python bindings}

All GPU functions are exposed via a pybind11 module
(\texttt{tmrca\_cu.\_core}).  Key entry points include:
\begin{itemize}
  \item \texttt{hmm\_posterior(G, positions, pair, ...)}: Single-pair
    full posterior (\textsc{Full\_Gamma}).
  \item \texttt{hmm\_posterior\_batched(G, positions, pairs, ...)}:
    Batched inference (\textsc{Summary\_Only}).
  \item \texttt{adaptive\_prior\_infer(G, positions, pairs, ...)}: EM
    loop (\textsc{Summary\_And\_EM}).
  \item \texttt{ep\_infer(G, positions, pairs, m, ...)}: EP loop with
    ultrametric constraints.
  \item \texttt{HMMContext(G, positions, ...)}: Persistent GPU context
    with auto-chunked \texttt{run\_batch()}.
\end{itemize}

A high-level \texttt{CoalescenceEstimator} Python class provides a
scikit-learn-style API with methods
\texttt{site\_pi()}, \texttt{pairwise\_divergence()}, and
\texttt{infer\_tmrca()}.

\subsection{Memory layout}

All arrays use row-major (C-order) layout.  Key strides:
\begin{align*}
  \texttt{packed}[i, w] &:\quad i \cdot n_\mathrm{words} + w,\\
  \texttt{gamma}[p, s, k] &:\quad (p \cdot S + s) \cdot K + k,\\
  \texttt{tmrca\_mean}[p, s] &:\quad p \cdot S + s.
\end{align*}


% ════════════════════════════════════════════════════════════════
\section{Validation}
\label{sec:validation}

\subsection{Reference implementation}
\label{sec:reference}

A pure-NumPy implementation (\texttt{NumpyHMM}) independently
implements the forward, backward, posterior, and log-likelihood
computations with identical mathematical formulas but in double
precision and without any GPU-specific optimizations.  This serves as
the ground truth for numerical validation.

The reference implementation uses explicit $\calO(K^2)$ loops for the
transition step (enumerating all source--target bin pairs), whereas the
GPU kernel exploits the SMC factorization to reduce this to
$\calO(K)$.  Both produce identical results up to floating-point
precision.

\subsection{Unit tests}
\label{sec:unit_tests}

The test suite (69~tests, \texttt{pytest}) covers:

\begin{itemize}
  \item \textbf{Bitpack roundtrip}: Exact recovery of genotypes
    after pack/unpack.
  \item \textbf{SFS}: Exact match with NumPy \texttt{bincount} and
    tskit \texttt{allele\_frequency\_spectrum}.
  \item \textbf{HMM posteriors}: GPU $\gamma_{s,k}$ vs.\ NumPy
    reference, $\texttt{rtol} = 10^{-2}$, $\texttt{atol} = 10^{-4}$.
  \item \textbf{Posterior normalization}: $\sum_k \gamma_s(k) = 1$
    at every site ($\texttt{rtol} = 10^{-3}$).
  \item \textbf{Log-likelihood}: GPU vs.\ NumPy,
    $\texttt{rtol} = 10^{-2}$.
  \item \textbf{Transition matrix}: Row sums $= 1$; identity at
    $r = 0$; stationary distribution is $q$.
  \item \textbf{Emission model}: Monotonicity, completeness, boundary
    behavior.
  \item \textbf{Numerical stability}: Long monomorphic stretches
    ($10$~Mb gaps), saturated divergence, extreme $\Ne$ values
    ($10^2$ to $10^7$).
  \item \textbf{Ultrametric projection}: Tree consistency,
    merge-time monotonicity, reproducibility.
\end{itemize}

\subsection{Statistical validation against msprime}
\label{sec:msprime_validation}

We simulate genealogies with known parameters using
msprime~\citep{Baumdicker2022} and compare inferred TMRCAs to the
true values extracted via \texttt{tree.tmrca(i,j)}.

\subsubsection{Demographic scenarios}
\label{sec:scenarios}

Four scenarios are used throughout validation and benchmarking
(Table~\ref{tab:scenarios}):

\begin{table}[h]
\centering
\caption{Demographic scenarios for validation.}
\label{tab:scenarios}
\begin{tabular}{ll}
\toprule
Scenario & Parameters \\
\midrule
Constant     & $\Ne = 10{,}000$ \\
Bottleneck   & $10{,}000 \to 500$ (at $t=2000$) $\to 20{,}000$ (at $t=3000$) \\
Exp.\ growth & $\Ne(0) = 100{,}000$; growth rate $= 0.01$ for $t < 200$ \\
Structured   & 2 populations, migration rate $m = 10^{-4}$ \\
\bottomrule
\end{tabular}
\end{table}

\subsubsection{Accuracy metrics}

For each scenario, we compute:
\begin{itemize}
  \item \textbf{Pearson $r$}: Per-site correlation between true and
    estimated TMRCA (target $r > 0.8$ for Tier~3).
  \item \textbf{RMSE}: Root mean squared error in generations.
  \item \textbf{95\% CI coverage}: Fraction of sites where the true
    TMRCA falls within the credible interval (target $\sim\!0.95$).
  \item \textbf{Prior recovery}: Comparison of the EM-estimated
    coalescent prior $\hat{q}$ with the true prior (Tier~4).
\end{itemize}


% ════════════════════════════════════════════════════════════════
\section{Results}
\label{sec:results}

\subsection{Kernel throughput}
\label{sec:throughput}

Table~\ref{tab:speedup} summarizes the cumulative effect of each
optimization phase, measured on an NVIDIA A100-SXM4 (80~GB,
\texttt{sm\_80}) with $K = 32$, $S = 50{,}000$~sites, and
varying numbers of pairs.

\begin{table}[h]
\centering
\caption{Cumulative kernel throughput after each optimization phase.
Measured throughput in millions of site$\cdot$pairs/s ($K = 32$,
$S = 50{,}000$).}
\label{tab:speedup}
\begin{tabular}{lccr}
\toprule
Phase & Individual & Cumulative & Msite$\cdot$pairs/s \\
\midrule
Baseline (FP64, na\"ive)      & ---               & $1.0\times$  & 9.1 \\
1.\ Emission precomp + \texttt{\_\_expf}
                               & $3\text{--}4\times$ & $3.5\times$  & $\sim\!32$ \\
2.\ FP32 throughout           & $\sim\!2\times$     & $\sim\!7\times$ & $\sim\!64$ \\
3.\ Multi-pair blocks (PPB=4) & $1.3\times$       & $\sim\!8\times$ & $\sim\!73$ \\
4.\ Fused summary (no D2H)    & $1.5\times$       & $\sim\!8.5\times$ & $\sim\!78$ \\
5.\ Persistent context        & $1.25\times$      & $10.2\times$ & 93 \\
\bottomrule
\end{tabular}
\end{table}

\noindent The throughput is consistent across scales: 92.3M at
4{,}950~pairs, 86.4M at 19{,}900~pairs (auto-chunked into 3~batches
of ${\sim}\,8{,}500$), and 92.8M at 124{,}750~pairs.

\begin{figure}[!ht]
\centering
\includegraphics[width=\textwidth]{fig3_speedup.pdf}
\caption{\textbf{Kernel optimization trajectory and throughput scaling.}
(\textit{Left})~Cumulative throughput (millions of site$\cdot$pairs/s)
after each of the five optimization phases, measured on an A100 with
$K = 32$, $S = 50{,}000$~sites.  The dominant contribution comes from
emission precomputation and fast-math intrinsics (Phase~1, $3.5\times$).
(\textit{Right})~Throughput as a function of the number of pairs,
demonstrating near-constant performance from 500 to 5{,}000~pairs.}
\label{fig:speedup}
\end{figure}

\subsection{Projected wall-clock times}
\label{sec:wallclock}

For human chromosome~1 ($S \approx 4.5 \times 10^6$ SNPs) at various
sample sizes:

\begin{table}[h]
\centering
\caption{Projected wall-clock time for all-pairs TMRCA inference on
human chr1 ($S = 4.5\text{M}$ SNPs), single A100.}
\label{tab:wallclock}
\begin{tabular}{rrrrr}
\toprule
$2N$ & Pairs & Site$\cdot$pairs & Baseline (9.1M/s) & Optimized (90M/s) \\
\midrule
100   & 4{,}950      & $2.2 \times 10^{10}$  & 40~min   & 4~min \\
500   & 124{,}750    & $5.6 \times 10^{11}$  & 17~hr    & 1.7~hr \\
1{,}000 & 499{,}500  & $2.25 \times 10^{12}$ & 69~hr    & \textbf{7~hr} \\
5{,}000 & 12.5M      & $5.6 \times 10^{13}$  & 71~days  & 7.2~days \\
\bottomrule
\end{tabular}
\end{table}

\begin{figure}[!ht]
\centering
\includegraphics[width=\textwidth]{fig6_wallclock.pdf}
\caption{\textbf{Projected wall-clock time for chromosome~1 all-pairs
inference.}  Log-log plot of total compute time versus sample size
($2N$) for human chr1 ($S = 4.5$M SNPs) on a single A100.  The
baseline kernel (9.1M site$\cdot$pairs/s, dashed) is compared with
the optimized kernel (93M site$\cdot$pairs/s, solid).  Horizontal
reference lines indicate 1~hour, 1~day, and 1~week.  At $2N = 1{,}000$,
the optimized kernel completes in ${\sim}\,7$~hours versus 69~hours at
baseline.}
\label{fig:wallclock}
\end{figure}

\subsection{Numerical accuracy}
\label{sec:accuracy}

\subsubsection{GPU versus NumPy reference}

The maximum element-wise relative difference between GPU and NumPy
posterior marginals is $< 10^{-2}$ across all tested configurations
($K \in \{32, 64, 128\}$, $S \in \{500, 50{,}000\}$,
$n \in \{20, 200\}$).  The \texttt{HMMContext} API produces bitwise
identical results to the standalone batched API (maximum difference
$= 0.0$), confirming that the persistent buffer path introduces no
numerical divergence.

\subsubsection{Correlation with true TMRCA}

Under the four demographic scenarios (Table~\ref{tab:scenarios}),
per-site Pearson correlations with true msprime coalescence times are:

\begin{table}[h]
\centering
\caption{Per-site Pearson $r$ between estimated and true TMRCA.
$n = 50$ diploid individuals, 1~Mb, $K = 64$.}
\label{tab:accuracy}
\begin{tabular}{lcc}
\toprule
Scenario & Tier~3 (const.\ $q$) & Tier~4 (adaptive $q$) \\
\midrule
Constant       & 0.80 & 0.85 \\
Bottleneck     & 0.75 & 0.82 \\
Exp.\ growth   & 0.72 & 0.80 \\
Structured     & 0.68 & 0.78 \\
\bottomrule
\end{tabular}
\end{table}

\noindent Tier~4 consistently outperforms Tier~3 by learning the
empirical coalescent time distribution, which is especially beneficial
under non-constant demography.

\begin{figure}[!ht]
\centering
\includegraphics[width=\textwidth]{fig2_accuracy.pdf}
\caption{\textbf{TMRCA estimation accuracy under constant demography.}
(\textit{Left})~Hexbin density of estimated versus true per-site TMRCA
for 45~pairs at 1~Mb ($r = 0.78$, $K = 32$, $\Ne = 10{,}000$).
(\textit{Centre})~Genomic landscape of estimated TMRCA (blue) with
95\% credible intervals (shaded band) versus true TMRCA (coral) for a
representative pair across a 200~kb window.
(\textit{Right})~Distribution of 95\% CI coverage across pairs
(median = 0.94).}
\label{fig:accuracy}
\end{figure}

\subsubsection{Credible interval calibration}

The 95\% credible intervals achieve coverage of
$0.90\text{--}0.95$ across scenarios, indicating well-calibrated
uncertainty quantification.  Under-coverage in the structured scenario
reflects the model mis-specification (constant $\Ne$ assumed, true
history involves migration).

\subsubsection{Prior recovery}

Under the bottleneck scenario, the EM-estimated prior after 10
iterations closely matches the true coalescent time distribution,
recovering the characteristic bimodal shape induced by the population
size change.  Under the constant scenario, convergence is achieved
in 3--5 iterations.

\begin{figure}[!ht]
\centering
\includegraphics[width=0.7\textwidth]{fig5_prior_recovery.pdf}
\caption{\textbf{EM recovery of the coalescent prior under a
bottleneck.}  Step plot comparing the true coalescent time
distribution (black) under a bottleneck scenario
($\Ne: 10{,}000 \to 500 \to 20{,}000$) with the initial constant-$\Ne$
prior (dashed grey) and the EM-estimated prior after 10~iterations
(blue, $K = 64$).  The EM successfully recovers the bimodal structure
characteristic of a population bottleneck.}
\label{fig:prior_recovery}
\end{figure}


% ════════════════════════════════════════════════════════════════
\section{Discussion}
\label{sec:discussion}

\subsection{Speedup analysis}

The $10.2\times$ speedup is attributable to five synergistic
optimizations.  The largest single contributor is emission
precomputation combined with fast-math intrinsics ($3\text{--}4\times$),
which eliminates the dominant cost of the baseline kernel: 8
double-precision \texttt{exp()} calls per site.  The FP32 conversion
($2\times$) exploits the A100's $2\!:\!1$ FP32/FP64 throughput ratio.
Multi-pair blocks improve SM occupancy but contribute modestly
($1.3\times$) because the kernel is not purely compute-bound.  The
fused summary extraction and persistent context phases each
contribute ${\sim}\,1.25\text{--}1.5\times$ by eliminating memory
transfer bottlenecks.

The observed speedup falls in the lower range of the initial estimate
(10--30$\times$).  The gap from the upper bound is explained by
memory bandwidth limitations: at $K = 32$, each site requires reading
two 64-bit packed genotype words, two \texttt{double} cumulative
recombination values, and two \texttt{double} position values, plus
writing \texttt{float} summary outputs.  The arithmetic intensity
(${\sim}\,100$~FLOP per ${\sim}\,40$~bytes read) places the kernel in
a regime where both compute and memory throughput matter.

\subsection{Scalability}

The auto-chunking mechanism in \texttt{HMMContext} enables transparent
scaling to arbitrary numbers of pairs without user-visible memory
management.  For the $2N = 1{,}000$ use case
($499{,}500$~pairs, $4.5$M~sites), the 80~GB A100 accommodates
${\sim}\,8{,}500$~pairs per chunk (limited by the alpha scratch buffer:
$8{,}500 \times 50{,}000 \times 32 \times 4\text{B} \approx 54$~GB).
The chunking overhead (pair index upload, result download) is
negligible relative to the kernel execution time.

Multi-GPU parallelism is straightforward: pairs are embarrassingly
parallel, so $p$~GPUs yield a near-linear $p\times$ speedup.  With
8~A100s, the chr1 $2N\!=\!1{,}000$ workload would complete in under
1~hour.

\subsection{Accuracy considerations}

The FP32 conversion introduces relative errors of ${\sim}\,10^{-3}$
per site, well within the statistical uncertainty of TMRCA estimates.
The per-site normalization prevents error accumulation across sites.
The \texttt{\_\_expf()} fast-math intrinsic has a worst-case error of
$2^{-21}$ (${\sim}\,5 \times 10^{-7}$), which is negligible compared
to the FP32 mantissa precision of $2^{-23}$.

The gap emission model assumes a uniform mutation rate across the
genome.  Incorporating a per-site rate map would require only minor
modifications (precomputing $K$ emission values per site, or reading
$\mu(s)$ from a device array at each site).

\subsection{Comparison with existing methods}

PSMC~\citep{Li2011} operates on a single diploid genome and requires
${\sim}\,1$~hour per sample on CPU\@.  MSMC2~\citep{Schiffels2020}
handles up to ${\sim}\,10$ haplotypes.
\texttt{Relate}~\citep{Speidel2019} scales to thousands of samples
but infers full ARGs rather than pairwise TMRCAs, with runtimes of
hours to days.

\texttt{tmrca.cu} occupies a different niche: it provides
site-resolved pairwise TMRCA posterior distributions (not point
estimates or tree topologies) at a throughput of $93$M
site$\cdot$pairs/s.  This enables analyses that require dense TMRCA
matrices, such as local IBD detection, TMRCA-based association
testing, and genome-wide genealogical distance computation.

\begin{figure}[!ht]
\centering
\includegraphics[width=\textwidth]{fig4_competitors.pdf}
\caption{\textbf{Benchmark against existing methods.}
(\textit{Left})~Single-pair inference time for sequences of 1--50~Mb:
PSMC (25~EM iterations), \texttt{gamma\_smc}, and \texttt{tmrca.cu}
on GPU\@.  \texttt{tmrca.cu} achieves 15--143$\times$ speedup over
\texttt{gamma\_smc} and orders of magnitude over PSMC\@.
(\textit{Right})~Multi-pair throughput at 1~Mb for 10, 50, and 100
haplotypes.  \texttt{tmrca.cu}'s batch parallelism yields
48--182$\times$ speedup, with the advantage growing as the number of
pairs increases.}
\label{fig:competitors}
\end{figure}

\subsubsection{Detailed scaling comparison with \texttt{gamma\_smc}}

We conduct a systematic comparison across a grid of sample sizes
($2$--$200$ haplotypes) and sequence lengths ($0.5$--$10$~Mb),
using msprime-simulated data under constant $\Ne = 10{,}000$
(Fig.~\ref{fig:gamma_deep}).

Three scaling regimes emerge.  First, at small sample sizes
($n = 2$--$10$), \texttt{gamma\_smc}'s ${\sim}\,6$s constant overhead
(VCF parsing and model initialization) dominates, yielding apparent
speedups of $200$--$450\times$.  Second, at moderate sample sizes
($n = 20$--$100$), both tools scale quadratically with the number of
haplotypes, but \texttt{tmrca.cu}'s GPU parallelism maintains a
$16$--$190\times$ advantage.  Third, at the largest configurations
($n = 200$, $10$~Mb), the speedup narrows to ${\sim}\,13\times$ as
\texttt{tmrca.cu} approaches GPU memory limits and must auto-chunk
into multiple batches.

The amortized per-pair cost (Fig.~\ref{fig:gamma_throughput}b)
reveals the underlying scaling: \texttt{gamma\_smc}'s per-pair time
decreases from ${\sim}\,6$~s (1~pair, dominated by I/O) to
${\sim}\,0.7$~ms ($19{,}900$~pairs), while \texttt{tmrca.cu}'s drops
from ${\sim}\,300$~ms (1~pair, kernel launch overhead) to
${\sim}\,0.04$~ms ($19{,}900$~pairs).  At saturation, the GPU
achieves ${\sim}\,77$M site$\cdot$pairs/s versus
${\sim}\,4$M for \texttt{gamma\_smc}, a $19\times$ throughput
ratio that reflects the fundamental hardware advantage.

\begin{figure}[!ht]
\centering
\includegraphics[width=\textwidth]{fig7_gamma_deep.pdf}
\caption{\textbf{Scaling comparison with \texttt{gamma\_smc}.}
(\textit{a})~Wall-clock time versus number of haplotypes at 1~Mb;
annotations show the speedup factor at each point.
(\textit{b})~Wall-clock time versus sequence length at $n = 50$
haplotypes.
(\textit{c})~Speedup heatmap across the full grid of sample sizes and
sequence lengths.  Speedup ranges from $13\times$ ($n = 200$, $5$~Mb)
to $453\times$ ($n = 10$, $1$~Mb).}
\label{fig:gamma_deep}
\end{figure}

\begin{figure}[!ht]
\centering
\includegraphics[width=0.85\textwidth]{fig8_gamma_throughput.pdf}
\caption{\textbf{Throughput and amortized per-pair cost at 1~Mb.}
(\textit{a})~Throughput in millions of site$\cdot$pairs/s for each
tool at varying sample sizes.  \texttt{tmrca.cu} maintains
$10$--$100\times$ higher throughput across all configurations.
(\textit{b})~Amortized per-pair inference time (ms) versus number of
pairs.  Both tools exhibit decreasing per-pair cost due to I/O
amortization, but \texttt{tmrca.cu}'s GPU parallelism yields a
steeper decrease, saturating at ${\sim}\,0.04$~ms/pair.}
\label{fig:gamma_throughput}
\end{figure}

\subsection{Limitations and future work}

\begin{enumerate}
  \item \textbf{Constant $\Ne$ assumption}: The prior
    (Eq.~\ref{eq:coal_prior}) assumes constant population size.
    While Tier~4 partially addresses this via EM, a fully
    parameterized piecewise-constant $\Ne(t)$ model would improve
    accuracy under complex demography.

  \item \textbf{Uniform mutation rate}: The current implementation
    uses a scalar $\mu$ across all sites.  Per-site rate maps (e.g.,
    from de novo mutation studies) could be incorporated with minimal
    kernel changes.

  \item \textbf{Single GPU}: The current implementation runs on a
    single device.  Multi-GPU and multi-node extensions via pair
    sharding are straightforward.

  \item \textbf{Phasing requirement}: The model assumes phased
    haplotypes.  Extending to unphased diploid data would require
    marginalizing over phase configurations.
\end{enumerate}


% ════════════════════════════════════════════════════════════════
\section{Conclusion}
\label{sec:conclusion}

We have presented \texttt{tmrca.cu}, a GPU-accelerated system for
pairwise TMRCA inference that achieves $93 \times 10^6$
site$\cdot$pairs/s on a single NVIDIA A100---a $10.2\times$ speedup
over the na\"ive baseline, making all-pairs analysis of
$2N\!=\!1{,}000$ across a human chromosome feasible in under 7~hours.
The four-tier architecture provides a spectrum of speed--accuracy
trade-offs, from instant windowed divergence to full posterior
inference with adaptive priors and ultrametric constraints.  Five
phases of systematic kernel optimization, together with fused summary
extraction and persistent GPU memory management, demonstrate how
careful co-design of numerical algorithms and GPU hardware can
unlock population-genomic analyses at biobank scale.


% ════════════════════════════════════════════════════════════════
\section*{Data and code availability}

\texttt{tmrca.cu} is implemented in CUDA~C++17 with pybind11 Python
bindings.  Source code, tests, and benchmarks are available at
\texttt{[repository URL]}.  Simulated data were generated using
msprime~v1.3~\citep{Baumdicker2022}.

\section*{Acknowledgments}

Computations were performed on an NVIDIA A100-SXM4 GPU\@.

% ════════════════════════════════════════════════════════════════
\bibliographystyle{plos2015}
\begin{thebibliography}{99}

\bibitem{Li2011}
Li, H.\ \& Durbin, R.
Inference of human population history from individual whole-genome sequences.
\textit{Nature} \textbf{475}, 493--496 (2011).

\bibitem{Schiffels2014}
Schiffels, S.\ \& Durbin, R.
Inferring human population size and separation history from multiple genome sequences.
\textit{Nat.\ Genet.} \textbf{46}, 919--925 (2014).

\bibitem{Schiffels2020}
Schiffels, S.\ \& Wang, K.
MSMC and MSMC2: The Multiple Sequentially Markovian Coalescent.
\textit{Methods Mol.\ Biol.} \textbf{2090}, 147--166 (2020).

\bibitem{Terhorst2017}
Terhorst, J., Kamm, J.A.\ \& Song, Y.S.
Robust and scalable inference of population history from hundreds of unphased whole genomes.
\textit{Nat.\ Genet.} \textbf{49}, 303--309 (2017).

\bibitem{Stern2019}
Stern, A.J., Wilton, P.R.\ \& Nielsen, R.
An approximate full-likelihood method for inferring selection and allele frequency trajectories from DNA sequence data.
\textit{PLoS Genet.} \textbf{15}, e1008384 (2019).

\bibitem{McVean2005}
McVean, G.A.T.\ \& Cardin, N.J.
Approximating the coalescent with recombination.
\textit{Philos.\ Trans.\ R.\ Soc.\ B} \textbf{360}, 1387--1393 (2005).

\bibitem{Kelleher2019}
Kelleher, J.\ et al.
Inferring whole-genome histories in large population datasets.
\textit{Nat.\ Genet.} \textbf{51}, 1330--1338 (2019).

\bibitem{Speidel2019}
Speidel, L., Forest, M., Shi, S.\ \& Myers, S.R.
A method for genome-wide genealogy estimation for thousands of samples.
\textit{Nat.\ Genet.} \textbf{51}, 1321--1329 (2019).

\bibitem{Zhang2023}
Zhang, B.C.\ et al.
Biobank-scale inference of ancestral recombination graphs.
\textit{Nat.\ Genet.} \textbf{55}, 768--776 (2023).

\bibitem{Killick2012}
Killick, R., Fearnhead, P.\ \& Eckley, I.A.
Optimal detection of changepoints with a linear computational cost.
\textit{J.\ Am.\ Stat.\ Assoc.} \textbf{107}, 1590--1598 (2012).

\bibitem{Baumdicker2022}
Baumdicker, F.\ et al.
Efficient ancestry and mutation simulation with msprime 1.0.
\textit{Genetics} \textbf{220}, iyab229 (2022).

\bibitem{prefix_pixi}
prefix.dev.\ pixi: Fast software package manager.
\url{https://pixi.sh} (2024).

\end{thebibliography}

% ════════════════════════════════════════════════════════════════
\appendix

\section{Detailed kernel pseudocode}
\label{app:pseudocode}

Algorithm~\ref{alg:forward_backward} provides the complete pseudocode
for the fused forward--backward kernel in \textsc{Summary\_Only} mode.

\noindent\textbf{Algorithm 1.} Batched HMM forward--backward kernel (one block, PPB pairs,
\textsc{Summary\_Only} mode).

\begin{Verbatim}[fontsize=\small,frame=single,label=Pseudocode]
Input:  packed genotypes, positions, mu, cum_rho, t_mid, q, pair indices
Output: T_hat[s], T_lo[s], T_hi[s], log_lik  (per pair)

k <- threadIdx.x mod K;  p <- threadIdx.x / K
(hi, hj) <- pair indices for pair p

// Phase 1: Precompute emissions
e_match    <- __expf(-2 * mu0 * t_mid[k])
e_mismatch <- 1 - e_match
g_base     <- -2 * mu0 * t_mid[k]

// Forward pass
alpha <- q[k] * emit(d[0], k);  normalize;  LL <- log(Z_0)
store[0, k] <- alpha

for s = 1 to S-1:
    r     <- cum_rho[s] - cum_rho[s-1]
    stay  <- __expf(-r * t_mid[k]) * alpha
    rm    <- (1 - __expf(-r * t_mid[k])) * alpha
    R     <- WarpReduce(rm)           // sum_l recomb mass
    alpha <- stay + q[k] * R
    alpha <- alpha * gap(s,k) * emit(d[s], k)
    normalize;  LL <- LL + log(Z_s)
    store[s, k] <- alpha

// Backward pass with fused summary extraction
beta <- 1
for s = S-2 downto 0:
    b[k]  <- beta * emit(d[s+1], k) * gap(s+1, k)
    Q     <- WarpReduce(q[k] * b[k])  // sum_l q_l * b_l
    beta  <- exp(-r * t_mid[k]) * b[k] + (1 - exp(-r * t_mid[k])) * Q
    normalize beta
    gamma <- store[s, k] * beta;  normalize
    T_hat[s] <- WarpReduce(gamma * t_mid[k])   // mean
    (T_lo[s], T_hi[s]) <- CdfScan(gamma)       // 95% CI via shuffles
    write T_hat[s], T_lo[s], T_hi[s]

write LL
\end{Verbatim}

\section{Memory budget}
\label{app:memory}

Table~\ref{tab:memory} details the GPU memory requirements for the
HMM inference at representative scales.

\begin{table}[h]
\centering
\caption{GPU memory budget per component.}
\label{tab:memory}
\begin{tabular}{lrl}
\toprule
Buffer & Size formula & Example ($n_p\!=\!8{,}500$, $S\!=\!50$K, $K\!=\!32$) \\
\midrule
Bitpacked genotypes & $n \cdot \lceil S/64 \rceil \cdot 8$ & 6.1~MB ($n\!=\!1{,}000$) \\
Alpha scratch       & $n_p \cdot S \cdot K \cdot 4$       & 54.4~GB \\
Summary arrays ($\times 3$) & $n_p \cdot S \cdot 4 \cdot 3$ & 5.1~GB \\
Log-likelihood      & $n_p \cdot 8$                        & 68~KB \\
Pair indices        & $n_p \cdot 4 \cdot 2$                & 68~KB \\
Parameters          & $(2S + 2K) \cdot 8$                  & 0.8~MB \\
\midrule
\textbf{Total}      &                                      & $\sim\!60$~GB \\
\bottomrule
\end{tabular}
\end{table}

\noindent The alpha scratch buffer dominates.  On an 80~GB A100 with
512~MB reserved, the auto-chunker computes
$n_\mathrm{chunk} = \lfloor (80 - 0.5 - 6.9) / 7.0 \rfloor \approx 10{,}500$
pairs per chunk (where 7.0~MB $\approx S \cdot K \cdot 4 + 3 \cdot S
\cdot 4 + 16$ per pair, and 6.9~GB covers static buffers).

\end{document}
